% =========================================================================
% Presentación científica del sistema de recuperación multisalto
% (grafo de aserciones reificadas + memoria canónica + planificador)
% y su aplicación al asistente clínico sobre las guías de GeSIDA.
%
% Compilar:  tectonic presentacion_sistema.tex
% =========================================================================
\documentclass[aspectratio=169,10pt]{beamer}
\usepackage[T1]{fontenc}
\usepackage[utf8]{inputenc}
\usepackage[spanish,es-noshorthands]{babel}
\usepackage{lmodern}
\usepackage{amsmath,amssymb}
\usepackage{booktabs}
\usepackage{tikz}
\usetikzlibrary{arrows.meta,positioning,fit,backgrounds,calc,shapes.geometric}
\usepackage{pgfplots}
\pgfplotsset{compat=1.17}

% ---------- estilo sobrio
\usetheme{default}
\usecolortheme{seagull}
\definecolor{azul}{RGB}{31,80,154}
\definecolor{teja}{RGB}{186,84,32}
\definecolor{verde}{RGB}{34,120,74}
\definecolor{grisf}{RGB}{100,100,100}
\setbeamercolor{frametitle}{fg=azul}
\setbeamercolor{title}{fg=azul}
\setbeamercolor{structure}{fg=azul}
\setbeamertemplate{navigation symbols}{}
\setbeamertemplate{footline}[frame number]
\setbeamertemplate{itemize items}[circle]
\newcommand{\sys}{\textsc{CatRAG-Reif}}
\newcommand{\hippo}{HippoRAG~2}
\newcommand{\bien}[1]{\textcolor{verde}{\textbf{#1}}}
\newcommand{\mal}[1]{\textcolor{teja}{\textbf{#1}}}

\tikzset{
  caja/.style={draw=azul, rounded corners=2pt, align=center, inner sep=4pt, fill=azul!8, font=\scriptsize},
  dato/.style={draw=grisf, dashed, rounded corners=2pt, align=center, inner sep=3pt, font=\scriptsize},
  ent/.style={draw=azul, ellipse, align=center, inner sep=2pt, font=\scriptsize},
  aser/.style={draw=azul, rectangle, rounded corners=1pt, align=center, inner sep=3pt, fill=azul!12, font=\scriptsize},
  chk/.style={draw=grisf, rectangle, align=center, inner sep=3pt, font=\scriptsize},
  fl/.style={-{Stealth[length=2mm]}, semithick, azul},
  flmal/.style={-{Stealth[length=2mm]}, semithick, teja},
  flok/.style={-{Stealth[length=2mm]}, semithick, verde},
  et/.style={font=\tiny, text=grisf, align=center},
}

\title{Recuperación multisalto sobre un grafo de aserciones reificadas\\
con memoria canónica de entidades y planificación de la evidencia}
\subtitle{Evaluación en benchmarks públicos y aplicación al asistente clínico
sobre las guías de GeSIDA}
\author{}
\date{Septiembre de 2026}

\begin{document}

\begin{frame}
\titlepage
\end{frame}

\begin{frame}{Guion}
\tableofcontents
\end{frame}

% =========================================================================
\section{Introducción}

\begin{frame}{El problema: la pregunta no se parece a la evidencia que falta}
\centering
\begin{tikzpicture}[node distance=4mm and 10mm]
  \node[dato, fill=azul!5, text width=44mm] (q) {\textbf{Pregunta}\\
    \emph{«Where did Sonia Greene's husband die?»}};
  \node[chk, right=18mm of q, text width=40mm] (p1) {\textbf{Pasaje 1} · \emph{Sonia Greene}\\
    «\dots{}her two-year marriage to \textbf{H.\,P. Lovecraft}\dots»};
  \node[chk, below=7mm of p1, text width=40mm] (p2) {\textbf{Pasaje 2 (puente)} · \emph{H. P. Lovecraft}\\
    «Howard Phillips Lovecraft\dots{} \textbf{died} in \textbf{Providence}\dots»};
  \draw[flok] (q) -- node[et, above]{$\cos=0{,}534$ · rango 1} (p1);
  \draw[flmal] (q.south east) to[out=-15,in=180]
      node[et, below, sloped]{$\cos=\mathbf{0{,}089}$ · rango \textbf{4\,358} de 6\,119} (p2);
  \node[et, below=1mm of p2, text width=52mm, align=center]
      {la respuesta está aquí, pero el vector de la pregunta\\ no contiene «Lovecraft»};
\end{tikzpicture}

\vspace{2mm}
\begin{itemize}
\item La recuperación densa/léxica encuentra lo \emph{nombrado}, no lo \emph{encadenado}:
  la información que identifica al pasaje puente solo aparece \emph{tras leer} el primero.
\item Consecuencia medible: exhaustividad parcial alta con \mal{cadena rota}
  (BM25 en cadenas dobles: $\mathrm{R@}5=50{,}9$ pero $\mathrm{FC@}5=0{,}9$).
\end{itemize}
\end{frame}

\begin{frame}{La métrica alineada con la tarea: cadena completa (FC@$k$)}
\begin{columns}[T]
\column{0.52\textwidth}
\begin{block}{Definiciones (pregunta $q$, oro $G_q$, top-$k$ $T_k$)}
\[
\mathrm{R@}k=\frac{|G_q\cap T_k|}{|G_q|}
\qquad
\mathrm{FC@}k=\mathbf{1}\!\left[G_q\subseteq T_k\right]
\]
\end{block}
\begin{itemize}
\item Un lector no puede componer la respuesta si le falta \emph{un} eslabón:
  $\mathrm{R@}5=0{,}75$ en una cadena de 4 vale lo mismo que 0.
\item FC@$k$ (\emph{Full Chain Retrieval}) es la métrica de integridad del
  razonamiento propuesta por CatRAG (Lau et al., ACL~2026).
\end{itemize}
\column{0.44\textwidth}
\centering
\begin{tikzpicture}[node distance=2.5mm]
  \node[chk, fill=verde!15] (a) {film A};
  \node[chk, fill=verde!15, right=of a] (b) {director A};
  \node[chk, fill=verde!15, below=of a] (c) {film B};
  \node[chk, fill=teja!20, right=of c] (d) {¿director B?};
  \draw[flok] (a) -- (b);
  \draw[flmal] (c) -- (d);
  \node[et, below=1mm of c, xshift=10mm, text width=40mm, align=center]
    {3 de 4 pasajes $\Rightarrow$ R@5 alta,\\ FC@5 $=0$: \mal{no hay respuesta fundamentada}};
\end{tikzpicture}
\end{columns}
\vspace{2mm}
\textbf{Estado del arte}: \hippo{} (PPR sobre grafo de tripletes; Gutiérrez et al.\ 2025)
y CatRAG (recorrido adaptativo a la consulta). Sus dos límites estructurales:
\emph{identidad de entidades sin resolver} y \emph{recursos fijos por consulta}.
\end{frame}

% =========================================================================
\section{Metodología}

\begin{frame}{Arquitectura: pagar una vez en indexación, decidir bien en consulta}
\centering
\begin{tikzpicture}[node distance=3mm and 4mm]
  % indexacion
  \node[dato] (corpus) {Corpus\\ $6\,119$ pasajes};
  \node[caja, right=of corpus] (openie) {Extracción\\ \textbf{reificada} (LLM)};
  \node[caja, right=of openie] (dicc) {\textbf{Memoria canónica}\\ de entidades};
  \node[caja, right=of dicc] (grafo) {Grafo canónico\\ $81$k nodos};
  \node[dato, right=of grafo] (emb) {vectores: hechos,\\ fichas, pasajes};
  \draw[fl] (corpus) -- (openie); \draw[fl] (openie) -- (dicc);
  \draw[fl] (dicc) -- (grafo); \draw[fl] (grafo) -- (emb);
  % consulta
  \node[dato, below=10mm of corpus] (q) {Pregunta};
  \node[caja, right=of q, fill=teja!12, draw=teja] (plan) {\textbf{Analista} (LLM)\\ menciones + \textbf{plan} + selección};
  \node[caja, right=of plan] (siembra) {Siembra\\ compuerta dura};
  \node[caja, right=of siembra] (ppr) {PPR\\ modulado};
  \node[caja, right=of ppr, fill=teja!12, draw=teja] (salto) {\textbf{Salto dirigido}\\ si quedan huecos};
  \node[caja, right=of salto, fill=teja!12, draw=teja] (conj) {\textbf{Conjunto}\\ 5 que cubren el plan};
  \draw[fl] (q) -- (plan); \draw[fl] (plan) -- (siembra);
  \draw[fl] (siembra) -- (ppr); \draw[fl] (ppr) -- (salto); \draw[fl] (salto) -- (conj);
  \draw[fl, grisf] (grafo.south) to[out=-90,in=90] (ppr.north);
  \begin{scope}[on background layer]
    \node[fit=(corpus)(emb), draw=grisf, rounded corners=3pt, inner sep=3pt,
      label={[et]above left:{\textbf{indexación} · una vez por corpus · sin ver ninguna pregunta}}] {};
    \node[fit=(q)(conj), draw=grisf, rounded corners=3pt, inner sep=3pt,
      label={[et]above left:{\textbf{consulta} · $\sim$2 llamadas LLM · $\sim$0{,}001\,\$}}] {};
  \end{scope}
\end{tikzpicture}

\vspace{2mm}
\begin{itemize}
\item Modelos pequeños en todo: \texttt{gpt-4o-mini} / \texttt{deepseek-chat}
  + \texttt{text-embedding-3-small} (1536 d).
\item En \textcolor{teja}{naranja}, las tres decisiones LLM de la consulta
  (el \emph{planificador}); el resto es álgebra y estructura.
\end{itemize}
\end{frame}

\begin{frame}{Pieza 1 · La unidad semántica: aserción reificada, no triplete}
\centering
\begin{tikzpicture}[node distance=4mm and 8mm]
  % izquierda: triple
  \node[ent] (t1) {Atomised};
  \node[ent, right=12mm of t1] (t2) {Oskar Roehler};
  \draw[fl] (t1) -- node[et, above]{\emph{directed}} (t2);
  \node[ent, below=5mm of t1, densely dotted] (t3) {Atomised (film)};
  \draw[grisf, densely dotted] (t1) -- node[et, left]{¿sinonimia?} (t3);
  \node[et, below=1mm of t3, xshift=8mm, text width=42mm, align=center]
    {\textbf{(a) \hippo{}}: el hecho es una \emph{arista}\\ binaria, anónima, sin procedencia};
  % derecha: asercion
  \node[aser, right=30mm of t2, text width=34mm] (a) {\textbf{aserción} $a$\\[-1pt]
    $\tau_a$: «Atomised was directed\\ by Oskar Roehler.»\\[1pt]
    {\tiny calificadores: fecha, lugar}};
  \node[ent, above left=4mm and -2mm of a] (e1) {can:atomised};
  \node[ent, above right=4mm and -2mm of a] (e2) {can:oskar roehler};
  \node[chk, below=5mm of a] (c1) {pasaje «Atomised (film)»};
  \draw[fl] (a) -- node[et, left]{\textsc{sujeto}} (e1);
  \draw[fl] (a) -- node[et, right]{\textsc{objeto}} (e2);
  \draw[fl] (a) -- node[et, right]{\textsc{procedencia}} (c1);
  \node[et, below=1mm of c1, text width=48mm, align=center]
    {\textbf{(b) \sys{}}: el hecho es un \emph{nodo} con frase propia,\\
     roles $n$-arios y vector $e(\tau_a)$};
\end{tikzpicture}

\vspace{1mm}
\begin{itemize}
\item Se incrusta \textbf{la frase}, no la etiqueta: robusto a etiquetas ruidosas;
  las fechas viajan \emph{dentro} del hecho.
\item El paseo aleatorio \textbf{razona a través del hecho}
  (entidad $\to$ aserción $\to$ otra entidad / pasaje) y su tránsito se modula
  por la consulta: $\mu_a=\tfrac14+\tfrac34\,\tilde\sigma_a$.
\end{itemize}
\end{frame}

\begin{frame}{Pieza 2 · Memoria canónica de entidades: matar las «islas» en el índice}
\begin{columns}[T]
\column{0.5\textwidth}
\centering
\begin{tikzpicture}[node distance=4mm]
  \node[ent, draw=teja] (m1) {«H.\,P. Lovecraft»\\ \tiny(en \emph{Sonia Greene})};
  \node[ent, draw=teja, below=of m1] (m2) {«Howard Phillips Lovecraft»\\ \tiny(en su biografía)};
  \draw[flmal, dashed] (m1) -- node[et, right]{$\cos = 0{,}781 < 0{,}80$\\ \mal{sin arista}} (m2);
  \node[et, below=1mm of m2, text width=50mm, align=center]
   {sinonimia por umbral (\hippo{}: $172$k aristas):\\ falla justo donde más se necesita\\
    \mal{57\,\% de los fallos de cadena eran islas}};
\end{tikzpicture}
\column{0.5\textwidth}
\centering
\begin{tikzpicture}[node distance=3mm]
  \node[aser, text width=52mm] (f) {\textbf{entrada canónica} \texttt{can:h.\,p.\,lovecraft}\\[-1pt]
   alias: \{H.\,P. Lovecraft; Howard Phillips Lovecraft\}\\[-1pt]
   descripción: «American writer of weird fiction\dots»\\[-1pt]
   \textbf{pasaje propio}: \emph{H. P. Lovecraft}};
  \node[et, below=1mm of f, text width=56mm, align=center]
   {fichas por mención (nombre $+$ contexto) $\to$ candidatos por 2 canales
    $\to$ restricciones duras (fechas, ordinales) $\to$ \textbf{adjudicación LLM}
    $\to$ fusión; el \textbf{título} del artículo es alias del sujeto};
\end{tikzpicture}
\end{columns}
\vspace{2mm}
\begin{itemize}
\item Construida \textbf{una vez, solo del corpus} (jamás del oro); auditable; con QIDs de
  Wikidata como verificación externa opcional (v4).
\item Resultado: variantes de nombre colapsan en \textbf{un} nodo $\Rightarrow$ el grafo
  canónico no necesita aristas de sinonimia; $178/188$ islas resueltas en el sondeo.
\item Es el análogo, derivado del corpus, del \emph{linker} SNOMED del caso clínico.
\end{itemize}
\end{frame}

\begin{frame}{Pieza 3 · El analista: planificar la cadena antes de seleccionar}
\footnotesize
\emph{«Which film has the director died first, Bílá Spona or A Night at Earl Carroll's?»}
\centering
\begin{tikzpicture}[node distance=2.5mm and 8mm]
  \node[dato, text width=51mm, align=left] (cand) {\textbf{hechos candidatos} (top-40 coseno $\cup$ frontera)\\
    2: A Night at Earl Carroll's \dots{} \textbf{directed by Kurt Neumann}\\
    6: Bílá spona was \dots{} \textbf{directed by Martin Fri\v{c}}\\
    8: \emph{Earl Carroll} (el actor) \emph{died on June 17, 1948}\, \dots};
  \node[caja, right=of cand, fill=teja!12, draw=teja, text width=52mm, align=left] (plan) {\textbf{plan} (huecos):\\
    director de Bílá spona $\to$ Martin Fri\v{c}\\
    fecha de muerte de Martin Fri\v{c}\\
    director de A Night\dots{} $\to$ Kurt Neumann\\
    fecha de muerte de Kurt Neumann};
  \draw[fl] (cand) -- (plan);
\end{tikzpicture}

\vspace{1.5mm}
\begin{columns}[T]
\column{0.55\textwidth}
\textbf{Sin plan} (selector directo, tope 4):\\
\mal{eligió al actor Earl Carroll como «director»}\\ y omitió a Kurt Neumann
\emph{que estaba en la lista} $\Rightarrow$ cadena rota.
\column{0.45\textwidth}
\textbf{Con plan} (tope $=$ n.º de huecos):\\
\bien{selecciona los dos hechos puente}; regla dura:
prohibido rellenar huecos con conocimiento propio del modelo ($[X]$).
\end{columns}
\vspace{1.5mm}
La misma llamada hace de \textbf{enlazador} (menciones $\to$ forma canónica $\to$ diccionario)
y declara los \textbf{huecos sin cubrir}, que alimentan el salto dirigido
\emph{y}, en clínica, la abstención fundamentada.
\end{frame}

\begin{frame}{Pieza 4 · Salto dirigido estructural: cuando el artículo habla en otro idioma}
\centering
\begin{tikzpicture}[node distance=3mm and 9mm]
  \node[dato, text width=40mm] (hueco) {\textbf{hueco del plan}\\ «director of \emph{The Magician} (1958 film)» $\to$ $[X]$};
  \node[caja, right=of hueco, text width=38mm] (dicc) {diccionario:\\ \texttt{can:the magician}\\ pasaje propio $=$\\ \emph{The Magician (1958 film)}};
  \node[dato, right=of dicc, text width=44mm] (hechos) {aserciones \textbf{de ese pasaje}:\\
    «\textbf{Ingmar Bergman} directed \textbf{Ansiktet}.»\\ \tiny(título original sueco)};
  \draw[fl] (hueco) -- node[et, above]{estructura,\\ no coseno} (dicc);
  \draw[fl] (dicc) -- (hechos);
  \node[et, below=2mm of dicc, text width=100mm, align=center]
   {el coseno del hueco con «Ingmar Bergman directed Ansiktet» es \mal{bajo} (cero vocabulario común);\\
    la ruta \bien{estructural} (entidad enlazada $\to$ su pasaje propio $\to$ sus hechos, con procedencia visible)\\
    lo rescata: el mini-filtro aprueba, Bergman se siembra, su biografía entra en el top-2};
\end{tikzpicture}
\vspace{2mm}
\begin{itemize}
\item Solo se ejecuta si el plan dejó huecos ($\sim$15--25\,\% de las consultas); una ronda.
\item Es la capacidad nueva de verdad: resolver puentes \textbf{inmunes al alias}
  (idiomas, títulos originales), imposibles para similitud y sinonimia.
\end{itemize}
\end{frame}

\begin{frame}{Piezas 5--6 · Siembra con compuerta dura y PPR modulado; conjunto final}
\begin{columns}[T]
\column{0.55\textwidth}
\textbf{Vector de reinicio} $\mathbf v$ (4 fuentes):
\begin{enumerate}\footnotesize
\item \textbf{compuerta dura}: solo siembran las entidades de hechos
  \emph{aprobados} (peso 1); si nada aprueba $\to$ degradación suave al denso;
\item presa densa sobre \emph{todos} los pasajes ($\times\,0{,}05$);
\item anclas débiles ($\epsilon=0{,}05$);
\item \textbf{pasaje propio}: $\mathbf v(\pi(\varepsilon)) \mathrel{+}= 0{,}5\,\mathbf v(\varepsilon)$
  --- «el artículo de una entidad \emph{es} la entidad».
\end{enumerate}
\[
\mathbf p^{(t+1)}=\lambda\!\left(P^{\top}\mathbf p^{(t)}+m_{\text{colg}}\mathbf v\right)+(1-\lambda)\,\mathbf v,
\quad \lambda=0{,}5
\]
con $P$ renormalizada tras modular las entradas a cada hecho por su afinidad
($\mu_a\in[\tfrac14,1]$; los aprobados a peso completo).
\column{0.45\textwidth}
\textbf{Conjunto final} (salvo comparaciones simples):\\[1mm]
\begin{tikzpicture}[node distance=2mm]
  \node[dato, text width=40mm] (top) {top-12 del PPR\\ \tiny(título + resumen)};
  \node[caja, below=of top, fill=teja!12, draw=teja, text width=40mm] (set)
    {LLM elige \textbf{el conjunto de 5}\\ que cubre \emph{todos} los huecos\\ \tiny(guarda: el top-2 del PPR se conserva)};
  \draw[fl] (top) -- (set);
\end{tikzpicture}\\[1mm]
\footnotesize FC@$k$ es una propiedad \emph{del conjunto}, no de cada pasaje:
esta decisión no la puede expresar un ranking punto a punto.
Aporta $+6$ puntos de FC@2 (orden temprano).
\end{columns}
\end{frame}

\begin{frame}{Protocolo experimental: paridad estricta y separación de datos}
\begin{columns}[T]
\column{0.52\textwidth}
\begin{itemize}\small
\item \textbf{Datasets} (subconjuntos de reproducción de \hippo{}/CatRAG):
  2WikiMultihopQA ($1\,000$ preguntas $\times$ $6\,119$ pasajes) y HotpotQA
  ($1\,000 \times 9\,811$); recuperación \emph{abierta} contra el corpus completo.
\item \textbf{Paridad}: \hippo{} ejecutado con \emph{su propio código} y los
  \emph{mismos modelos}; CatRAG publica con este mismo backbone
  (\texttt{gpt-4o-mini}+\texttt{te3-small}) $\Rightarrow$ comparación directa legítima.
\item \textbf{Validación cruzada de la reproducción}: nuestro \hippo{} en 2Wiki
  ($85{,}9/65{,}8$) coincide \bien{al decimal} con la reproducción independiente
  del artículo de CatRAG ($85{,}9/66{,}1$).
\item \textbf{Estadística}: pareado por pregunta, bootstrap $B{=}2\,000$, IC 95\,\%.
\end{itemize}
\column{0.46\textwidth}
\centering
\begin{tikzpicture}[node distance=3mm]
  \node[caja, fill=verde!12, draw=verde, text width=48mm] (cal)
    {\textbf{calibración} ($11\,376$ preguntas)\\ \tiny enrutador, ejemplos de consignas};
  \node[caja, fill=azul!10, below=of cal, text width=48mm] (son)
    {\textbf{sondeo} ($200$, retenido)\\ \tiny TODO el diseño y las ablaciones se deciden aquí};
  \node[caja, fill=teja!12, draw=teja, below=of son, text width=48mm] (banco)
    {\textbf{banco de pruebas} ($1\,000$)\\ \tiny intocable: una medición por configuración};
  \draw[fl] (cal) -- (son); \draw[fl] (son) -- (banco);
\end{tikzpicture}\\[1mm]
\footnotesize Configuración \textbf{congelada} antes de medir; en HotpotQA,
\emph{cero ajuste} al conjunto.
\end{columns}
\end{frame}

% =========================================================================
\section{Resultados}

\begin{frame}{2WikiMultihopQA: resultado principal (todos con el mismo backbone)}
\centering
\small
\begin{tabular}{lcccc}
\toprule
 & R@2 & R@5 & FC@5 & FC@10\\
\midrule
BM25 & --- & $65{,}8$ & $32{,}8$ & $40{,}1$\\
\hippo{} (su código) & $70{,}7$ & $85{,}9$ & $65{,}8$ & $71{,}7$\\
CatRAG (publicado; código no liberado) & --- & $87{,}0$ & $67{,}6$ & ---\\
\sys{} v3 (diccionario) & $73{,}0$ & $96{,}0$ & $90{,}6$ & $93{,}4$\\
\textbf{\sys{} v4 $+$ planificador} (\texttt{gpt-4o-mini}) & $80{,}2$ & $98{,}1$ & $96{,}0$ & $96{,}7$\\
\textbf{\sys{} v4 $+$ planificador} (\texttt{deepseek-chat}) & $\mathbf{84{,}2}$ & $\mathbf{98{,}9}$ & $\mathbf{97{,}7}$ & $\mathbf{97{,}8}$\\
\bottomrule
\end{tabular}

\vspace{2mm}
\begin{tikzpicture}
\begin{axis}[ybar, width=0.86\textwidth, height=38mm, bar width=7mm,
  symbolic x coords={BM25, HippoRAG 2, CatRAG (pub.), v3, v4 (4o-mini), v4 (deepseek)},
  xtick=data, ymin=0, ymax=105, ylabel={\scriptsize FC@5}, nodes near coords,
  every node near coord/.append style={font=\tiny},
  tick label style={font=\tiny}, label style={font=\tiny}]
\addplot[fill=azul!35, draw=azul] coordinates {(BM25,32.8) (HippoRAG 2,65.8) (CatRAG (pub.),67.6) (v3,90.6) (v4 (4o-mini),96.0) (v4 (deepseek),97.7)};
\end{axis}
\end{tikzpicture}

\scriptsize Pareado v4$-$\hippo{}: $\Delta\mathrm{FC@}5=+31{,}9$ $[+28{,}9,+34{,}9]$
(gana 326 / pierde 7); todos los tipos significativos. El artículo de \hippo{}
con NV-Embed-v2 (7\,B, GPU) reporta $\mathrm{R@}5=90{,}4$: aquí $98{,}9$ con
modelos de una fracción de ese coste.
\end{frame}

\begin{frame}{El hallazgo central: acumulación frente a colapso con la profundidad}
\begin{columns}[T]
\column{0.55\textwidth}
\centering
\begin{tikzpicture}
\begin{axis}[ybar, width=\textwidth, height=44mm, bar width=5.5mm,
  symbolic x coords={A sin puente, B puente simple, C doble puente},
  xtick=data, ymin=0, ymax=112, ylabel={\scriptsize FC@5},
  legend style={font=\tiny, at={(0.5,1.02)}, anchor=south, legend columns=2},
  nodes near coords, every node near coord/.append style={font=\tiny},
  tick label style={font=\tiny}, label style={font=\tiny}]
\addplot[fill=grisf!40, draw=grisf] coordinates {(A sin puente,93.4) (B puente simple,77.0) (C doble puente,12.3)};
\addplot[fill=verde!45, draw=verde] coordinates {(A sin puente,99.2) (B puente simple,89.6) (C doble puente,83.8)};
\legend{\hippo{}, \sys{} v3}
\end{axis}
\end{tikzpicture}
\column{0.45\textwidth}
\textbf{Modelo nulo}: eslabones independientes
$\Rightarrow \mathrm{FC}\approx p^{\,|G_q|}$
\begin{center}\small
\begin{tabular}{lccc}
\toprule
 & C predicho & C obs. & razón\\
\midrule
\hippo{} & $59{,}2$ & $12{,}3$ & \mal{0{,}21}\\
\sys{} & $80{,}3$ & $83{,}8$ & \bien{1{,}04}\\
\bottomrule
\end{tabular}
\end{center}
\footnotesize
\sys{} degrada por \emph{acumulación} de $\sim$5\,\% de error por eslabón.
\hippo{} \emph{colapsa}: sus recursos por consulta (selección $\le 4$, masa,
holgura del top-5) son fijos y compartidos entre las dos cadenas.
El planificador invierte el gradiente: con v4, la cadena doble alcanza
$\mathrm{FC@}5=98{,}7$ --- el segmento históricamente peor pasa a ser el mejor.
\end{columns}
\end{frame}

\begin{frame}{Atribución: ablaciones y control sin grafo}
\begin{columns}[T]
\column{0.5\textwidth}
\textbf{Escalera de contribuciones} (sondeo):
\begin{center}
\begin{tikzpicture}
\begin{axis}[width=\textwidth, height=40mm, ymin=88, ymax=100,
  symbolic x coords={v3, +plan, +salto, +conjunto},
  xtick=data, ylabel={\scriptsize FC@5}, nodes near coords,
  every node near coord/.append style={font=\tiny},
  tick label style={font=\tiny}, label style={font=\tiny}]
\addplot[mark=*, azul, thick] coordinates {(v3,92.0) (+plan,92.5) (+salto,95.5) (+conjunto,97.5)};
\end{axis}
\end{tikzpicture}
\end{center}
\footnotesize El plan por sí solo apenas mueve ($+0{,}5$): su valor es
\emph{habilitar} al salto ($+3{,}0$; cadena doble $86\to96$) y al conjunto
($+1{,}5$; $+6$ FC@2).
\column{0.5\textwidth}
\textbf{¿Grafo o planificador?} Control \emph{sin grafo} con idéntico
presupuesto LLM (mismo analista/salto/conjunto sobre top-20 denso):
\begin{center}\small
\begin{tabular}{lccc}
\toprule
banco & v4 & control & $\Delta$ [IC 95\,\%]\\
\midrule
2Wiki & $96{,}0$ & $83{,}3$ & \bien{$+12{,}7$} $[+10{,}5,+14{,}9]$\\
HotpotQA & $90{,}2$ & $83{,}9$ & \bien{$+6{,}3$} $[+4{,}1,+8{,}6]$\\
\bottomrule
\end{tabular}
\end{center}
\footnotesize Criterio pre-registrado ($\ge 5$ y sig.\ en ambos): \bien{cumplido}
$\Rightarrow$ la \textbf{representación aporta} más allá del planificador.
El planificador solo repara cadenas dobles; el puente simple lo repara el grafo
(oro perdido por el control: $130/141$ ni siquiera en su top-20).
\end{columns}
\end{frame}

\begin{frame}{Generalización sin ajuste y robustez al modelo}
\begin{columns}[T]
\column{0.55\textwidth}
\textbf{HotpotQA} ($1\,000\times9\,811$), configuración \emph{congelada} de 2Wiki:
\begin{center}\small
\begin{tabular}{lcc}
\toprule
 & R@5 & FC@5\\
\midrule
BM25 (nuestro) & $72{,}8$ & $49{,}5$\\
denso \texttt{te3-small} (pub.) & $81{,}3$ & $64{,}9$\\
\hippo{} (publicado) & $87{,}1$ & $75{,}5$\\
CatRAG (publicado) & $89{,}5$ & $80{,}4$\\
\textbf{\sys{} v4} (\texttt{4o-mini}) & $94{,}7$ & $90{,}2$\\
\textbf{\sys{} v4} (\texttt{deepseek}) & $\mathbf{96{,}0}$ & $\mathbf{93{,}0}$\\
\bottomrule
\end{tabular}
\end{center}
\footnotesize $+6{,}5$ R@5 / $+12{,}6$ FC@5 sobre CatRAG con \emph{cero} ajuste
al conjunto.
\column{0.45\textwidth}
\textbf{El modelo del planificador es intercambiable} (sondeo 2Wiki, FC@5):
\begin{center}\small
\begin{tabular}{lcc}
\toprule
modelo & FC@5 & \$/consulta\\
\midrule
deepseek-chat & $97{,}5$ & $0{,}0010$\\
groq gpt-oss-20b & $97{,}0$ & $0{,}0005$\\
gpt-4o & $96{,}5$ & $0{,}012$\\
gpt-4o-mini & $96{,}0$ & $0{,}0008$\\
\bottomrule
\end{tabular}
\end{center}
\footnotesize La ganancia es del \emph{diseño de la tarea}, no del tamaño:
un 20\,B abierto casi empata con \texttt{gpt-4o} a $1/24$ del coste.
Estudio completo: $\sim$40\,\$; toda llamada cacheada (reproducción a coste cero).
\end{columns}
\end{frame}

% =========================================================================
\section{Discusión}

\begin{frame}{Discusión: por qué funciona y qué falta}
\textbf{Tres principios, cada uno elimina una dependencia de la profundidad:}
\begin{enumerate}
\item \textbf{Identidad en el índice} (memoria canónica): las «islas» no se
  multiplican con los saltos; la sinonimia por umbral sí.
\item \textbf{Hechos como nodos de primera clase}: la masa viaja \emph{a través}
  del hecho aprobado; los calificadores viajan con él.
\item \textbf{Recursos que escalan con la cadena} (plan, salto, conjunto):
  el presupuesto se asigna por hueco, no por consulta --- de ahí la inversión
  del gradiente en cadenas dobles.
\end{enumerate}
\vspace{1mm}
\textbf{Validez}: reproducción del baseline verificada al decimal contra una
reproducción independiente; diseño íntegro en conjuntos retenidos; una medición
por configuración; artefactos y cachés públicos (reproducible a coste cero).

\textbf{Límites y trabajo en curso}: QA extremo a extremo (EM/F1, JSR) pendiente;
MuSiQue (saltos 2/3/4 etiquetados: contraste directo del modelo multiplicativo);
inferencia con $n{=}108$ no alcanza significación aislada; heurísticas léxicas
residuales sustituidas por el analista, pendiente de medir en otros idiomas.
\end{frame}

% =========================================================================
\section{Aplicación: asistente clínico sobre las guías de GeSIDA}

\begin{frame}{El caso de uso origen: asistente sobre 8 guías clínicas de GeSIDA}
\begin{columns}[T]
\column{0.55\textwidth}
La arquitectura no nació en Wikipedia: es la generalización del asistente
clínico VIH \textbf{ya construido}, donde cada pieza tiene su análogo curado:
\begin{center}\footnotesize
\begin{tabular}{ll}
\toprule
benchmark & clínica (GeSIDA)\\
\midrule
memoria canónica del corpus & \textbf{SNOMED\,CT/UMLS} (SCTID)\\
aserciones reificadas & aserciones con \textbf{deóntica},\\
 & grado de evidencia y \textbf{vigencia}\\
enrutador de intención & router clínico (paciente,\\
 & numérica) $+$ lente de umbrales\\
pasaje propio & sección normativa de la guía\\
\bottomrule
\end{tabular}
\end{center}
\column{0.45\textwidth}
\footnotesize
\textbf{Estado actual (producción, congelada)}:
\begin{itemize}\footnotesize
\item 8 guías; banco de $505$ preguntas con validación clínica en curso;
\item preguntas multi-guía (nivel 4): 84,9\,\% en el top-20 en modo grafo
  (denso: $24{,}7$; BM25: $45{,}2$) --- el análogo clínico del multisalto;
\item curación con humano en el bucle: bandeja de anclajes SNOMED auditable.
\end{itemize}
\end{columns}
\vspace{1mm}
\centering
\begin{tikzpicture}[node distance=2.5mm and 6mm]
  \node[dato, text width=40mm] (qc) {«¿Pauta de TAR en gestante con tuberculosis activa?»};
  \node[caja, right=of qc, fill=teja!12, draw=teja, text width=40mm] (planc) {\textbf{plan clínico}:\\ TAR en embarazo $\to$ guía Embarazo\\ interacción rifampicina $\to$ guía TB-VIH};
  \node[caja, right=of planc, text width=34mm] (ev) {cadena de evidencia\\ multi-guía $+$ vigencias};
  \draw[fl] (qc) -- (planc); \draw[fl] (planc) -- (ev);
\end{tikzpicture}
\end{frame}


\begin{frame}{GeSIDA de punta a punta: la metodología sobre un caso multi-guía}
\centering
\begin{tikzpicture}[node distance=1.6mm and 2.5mm, font=\tiny]
  % ---------- fila 1: pregunta -> analista -> linker
  \node[dato, text width=31mm] (qq) {\textbf{Pregunta}\\ «¿Se mantiene \textbf{dolutegravir} en una \textbf{gestante} que inicia \textbf{rifampicina}?»};
  \node[caja, right=of qq, fill=teja!12, draw=teja, text width=44mm, align=left] (an) {\textbf{Analista} · plan clínico (huecos):\\
    1. pauta de DTG en embarazo $\to$ guía \textbf{Embarazo}\\
    2. interacción DTG--rifampicina $\to$ guía \textbf{TB-VIH}\\
    3. ajuste de dosis si coadministración $\to$ $[X]$};
  \node[caja, right=of an, text width=33mm] (lk) {\textbf{Enlazador SNOMED}\\ dolutegravir $\to$ \texttt{sct:DTG} \checkmark\\ rifampicina $\to$ \texttt{sct:RIF} \checkmark\\ gestación $\to$ \texttt{sct:EMB} \checkmark};
  \draw[fl] (qq) -- (an); \draw[fl] (an) -- (lk);
  % ---------- fila 2: el grafo clinico (multisalto)
  \node[chk, below=7mm of qq, xshift=2mm, text width=27mm] (g1) {\textbf{guía Embarazo} · §pauta\\ \emph{pasaje propio} de la\\ recomendación en gestación};
  \node[aser, right=of g1, text width=36mm, align=left] (a1) {\textbf{aserción} $a_1$ · \textsc{RecomendacionTerapeutica}\\
    intervención: dolutegravir\\ población: \textbf{gestantes}\\ deóntica: SE RECOMIENDA (A-I)\\ vigencia: 2022 \checkmark};
  \node[ent, right=of a1, fill=verde!12] (dtg) {\texttt{sct:DTG}\\ \textbf{entidad puente}};
  \node[aser, right=of dtg, text width=37mm, align=left] (a2) {\textbf{aserción} $a_2$ · \textsc{Interaccion}\\
    fármacos: DTG $+$ \textbf{rifampicina}\\ efecto: inducción $\downarrow$ niveles\\ ajuste: \textbf{50 mg/12 h} (numérico)\\ deóntica: AJUSTAR DOSIS};
  \node[chk, right=of a2, text width=24mm] (g2) {\textbf{guía TB-VIH}\\ §interacciones};
  \draw[fl] (a1) -- (g1); \draw[fl] (dtg) -- (a1); \draw[fl] (dtg) -- (a2); \draw[fl] (a2) -- (g2);
  \node[et, below=0.5mm of dtg, text width=30mm]{siembra: compuerta dura\\ $+$ pasaje propio};
  % ---------- fila 3: ppr/salto -> conjunto -> respuesta
  \node[caja, below=8mm of g1, xshift=6mm, text width=33mm] (ppr2) {\textbf{PPR modulado}\\ la masa cruza $a_1$ y $a_2$ ($\mu{=}1$):\\ \emph{un solo grafo, dos guías}};
  \node[caja, right=of ppr2, fill=teja!12, draw=teja, text width=36mm] (sal) {\textbf{Salto dirigido} (hueco 3)\\ «dosis en coadministración» $\to$\\ aserciones de §interacciones\\ $\to$ aprueba $a_2$ \checkmark};
  \node[caja, right=of sal, fill=teja!12, draw=teja, text width=30mm] (cj) {\textbf{Conjunto final}\\ pasajes que cubren\\ los 3 huecos del plan};
  \node[dato, right=of cj, text width=27mm, fill=verde!8] (resp) {\textbf{Respuesta con expediente}\\ cadena $+$ procedencia\\ $+$ vigencia $+$ dosis\\[1pt]
    \mal{hueco sin cubrir} $\Rightarrow$\\ \textbf{abstención fundamentada}};
  \draw[fl] (ppr2) -- (sal); \draw[fl] (sal) -- (cj); \draw[fl] (cj) -- (resp);
  \draw[fl, grisf] (lk.south) to[out=-90,in=45] (dtg.north east);
  \draw[fl, grisf] (g1.south) to[out=-90,in=170] (ppr2.west);
\end{tikzpicture}

\vspace{1mm}
\scriptsize El multisalto clínico: la entidad canónica \texttt{sct:DTG} es el
\emph{puente} entre dos guías distintas --- exactamente la estructura del
puente de 2Wiki, con SNOMED en el papel de la memoria canónica. Contenido
abreviado con fines ilustrativos; las aserciones reales llevan guía, sección,
página y fecha de vigencia exactas.
\end{frame}

\begin{frame}{Qué aporta esta investigación a la clínica, y con qué salvaguardas}
\begin{columns}[T]
\column{0.55\textwidth}
\textbf{Los tres portes con evidencia detrás:}
\begin{enumerate}\small
\item \textbf{Cadena de evidencia multi-guía}: las preguntas nivel~4 son el
  análogo del puente; el gradiente invertido predice la mayor ganancia ahí.
\item \textbf{Abstención fundamentada}: los \emph{huecos sin cubrir} del plan
  se convierten en «no hay evidencia sobre $X$ en las guías» --- el fallo
  silencioso ($\mathrm{FC}=0$ respondido de memoria) pasa a ser un aviso
  explícito y medible por pregunta.
\item \textbf{JSR clínico}: éxito $=$ cadena normativa completa \emph{y}
  respuesta correcta \emph{y} vigencia comprobada --- respuesta con expediente.
\end{enumerate}
\column{0.45\textwidth}
\textbf{Salvaguardas no negociables:}
\begin{itemize}\small
\item Datos clínicos \textbf{nunca} a proveedores sin garantías (DeepSeek
  excluido del caso clínico); el diseño es agnóstico al modelo $\Rightarrow$
  proveedor UE o modelo local;
\item validación \emph{offline} contra el banco de 505 antes de tocar
  producción (sistema actual congelado y etiquetado);
\item terminología curada (SNOMED) como fuente de verdad de identidad;
  el LLM propone, \textbf{nunca decide identidades};
\item trazabilidad completa: cada respuesta con su cadena, procedencia
  y vigencia.
\end{itemize}
\end{columns}
\end{frame}

% =========================================================================
\begin{frame}{Conclusiones}
\begin{enumerate}
\item Un grafo de \textbf{aserciones reificadas} sobre \textbf{entidades
  canónicas}, consultado con un \textbf{planificador de la evidencia}, recupera
  la cadena completa donde el estado del arte la rompe:
  2Wiki $\mathrm{FC@}5$ $\mathbf{97{,}7}$ frente a $65{,}8$ (\hippo{}) y $67{,}6$
  (CatRAG), mismo backbone; HotpotQA $\mathbf{93{,}0}$ frente a $80{,}4$ con cero ajuste.
\item La mejora está \emph{atribuida}: representación ($+12{,}7$/$+6{,}3$ FC@5
  frente al control sin grafo con idéntico presupuesto LLM) $+$ planificador
  (repara las cadenas dobles: gradiente invertido) --- y es \emph{acumulación},
  no colapso ($\mathrm{FC}\approx p^{|G_q|}$).
\item Todo con modelos pequeños e intercambiables ($\sim$0{,}001\,\$/consulta),
  reproducible a coste cero, y con un camino directo a la clínica: mismas
  piezas con terminología curada, y la abstención fundamentada como dividendo
  de seguridad.
\end{enumerate}
\vspace{2mm}
\centering
\emph{Gracias --- preguntas.}\\[1mm]
\scriptsize código y artefactos: \texttt{github.com/bjur27r/hiv\_rag} (rama \texttt{eval-wiki2})
\end{frame}

\end{document}
